\section{第二层：物理与系统类比}

\subsection{为什么需要类比，但类比必须受约束}

生物学为本章提供正当性，但并不提供全部概念机器。
复杂系统与物理类比之所以有用，是因为它们为重复互动、协同锁定、阈值效应与多尺度秩序提供了有组织的思考方式。
它们在本章中的角色是\textbf{中等}强度且明确属于解释性的。

\begin{remark}
手稿并不声称书桌、研究实验室与神经元集合是同一种对象。
它所主张的是：它们可能共享一些可分析的性质，例如耦合、强化、路径依赖以及被重构后的成本景观。
\end{remark}

\subsection{自组织与序参量}

经典复杂性教材在这里之所以有用，是因为它们解释了一个基础思想：
局部互动可以产生高层协调，而不需要一个中央执行者预先显式计算完整模式。

用本手稿的语言来说，一个框架可以被看作许多局部规律的某种“序参量式”摘要：
\begin{itemize}[nosep]
  \item 一天中的时间；
  \item 线索摆放；
  \item 姿势；
  \item 奖赏预期；
  \item 社会同步；
  \item 对先前成功与先前失败的记忆。
\end{itemize}

一旦这些变量足够频繁地对齐，系统就不必每次都从零开始协商完整的行为选择。
一个宏观模式已经存在，而这个模式会约束下一轮微观决策。

\subsection{路径依赖与干扰}

路径依赖对于用户提出的“既往数据干扰”要求至关重要。
同样的当前刺激之所以可能导向不同动作，
是因为系统携带着不同的既往条件作用痕迹、既往损失、既往奖赏或既往环境配对。

\begin{definition}[路径依赖]
\textbf{路径依赖}是指当前转变成本依赖于先前状态与强化的序列，而不仅仅依赖当前快照。
\end{definition}

\begin{definition}[干扰]
\textbf{干扰}是指：先前被强化的路径、线索或预期，在当前情境中仍然保持活跃，从而扭曲或阻断目标转变。
\end{definition}

\begin{discussion}
这是本章在业务上最关键的观点之一。
坏框架之所以不断重建，往往不是因为当前决策在道德上“错误”，
而是因为旧路径仍然储存着更多结构。
同样，好框架之所以越来越容易，也不是因为行动者抽象意义上“变好了”，
而是因为这一条路径通过重复使用而变得更便宜。
\end{discussion}

\subsection{多尺度组织}

涌现之所以重要，还因为组织会同时存在于多个尺度上。
一个有用框架可以：
\begin{enumerate}[nosep]
  \item 在局部层面由小而重复的动作来实施；
  \item 由房间布局、日程结构或社会时序来维持；
  \item 嵌入在更大的生物学与意识状态条件之中。
\end{enumerate}

这一多尺度视角避免了一个常见错误：
试图仅用内在心理，或者仅用外部环境来解释全部框架形成。
本书的主张要求更高：只有当多个尺度对齐得足够久，一个路径才会相对其他路径稳定下来。